\documentclass[12pt,a4paper]{article}
\usepackage[utf8]{inputenc}
\usepackage[spanish]{babel}
\usepackage{amsmath}
\usepackage{amssymb}
\usepackage{geometry}
\usepackage{enumitem}
\usepackage{titlesec}
\usepackage{fancyhdr}
\usepackage{booktabs}
\usepackage{array}
\usepackage{mdframed}
\usepackage{xcolor}

\geometry{
  top=2.5cm,
  bottom=2.5cm,
  left=2.5cm,
  right=2.5cm
}

% Header and footer
\pagestyle{fancy}
\fancyhf{}
\rhead{DIS8012 -- 2026-1}
\lhead{Ayudantía 2}
\rfoot{Página \thepage}
\lfoot{Administración Financiera}

% Section formatting
\titleformat{\section}[block]
  {\large\bfseries\filcenter}
  {}
  {0em}
  {\colorbox{black!85}{\parbox{\dimexpr\textwidth-2\fboxsep\relax}{\centering\color{white}#1}}}
\titlespacing*{\section}{0pt}{1.2em}{0.8em}

\titleformat{\subsection}{\normalsize\bfseries}{}{0em}{}
\titlespacing*{\subsection}{0pt}{0.8em}{0.3em}

% Answer space commands
\newcommand{\answerspace}{\vspace{4.5cm}}
\newcommand{\lastanswerspace}{\vspace{2cm}}

% Formula box
\newmdenv[
  linecolor=black,
  linewidth=1pt,
  innertopmargin=6pt,
  innerbottommargin=6pt,
  innerleftmargin=10pt,
  innerrightmargin=10pt,
  backgroundcolor=gray!10
]{formulabox}

\begin{document}

% ─── TITLE ────────────────────────────────────────────────────────────────────
\begin{center}
  {\LARGE \textbf{Ayudantía 2 -- Guía de Ejercicios}}\\[0.4em]
  {\large DIS8012 · Administración Financiera · 2026-1}\\[0.2em]
  {\normalsize Resuelve cada ejercicio en el espacio indicado.}
\end{center}

\vspace{0.5em}
\noindent\rule{\textwidth}{0.8pt}
\vspace{0.3em}

% ─── FORMULARIO DE REFERENCIA ────────────────────────────────────────────────
\subsection*{Formulario de Referencia}

\begin{formulabox}
\begin{center}
\renewcommand{\arraystretch}{2.2}
\begin{tabular}{ll}
  \textbf{Costo Variable Unitario} & $cv = \dfrac{\text{Costos Variables Totales}}{\text{Unidades producidas}}$ \\[4pt]
  \textbf{Costo Fijo Unitario}     & $cf = \dfrac{\text{Costos Fijos Totales}}{\text{Unidades producidas}}$ \\[4pt]
  \textbf{Costo Total}             & $CT = F + (cv \times Q)$ \\[4pt]
  \textbf{Punto de Equilibrio}     & $Q_{eq} = \dfrac{F}{p - c}$ \\[4pt]
  \textbf{Ingresos Totales}        & $IT = p \times Q$ \\[4pt]
  \textbf{Valor Presente}          & $VP = \dfrac{VF}{(1+i)^n}$ \quad\quad \textbf{Valor Futuro:} $VF = VP \cdot (1+i)^n$ \\
\end{tabular}
\end{center}
\end{formulabox}

\vspace{0.6em}
\noindent\rule{\textwidth}{0.8pt}

% ═══════════════════════════════════════════════════════════════════════════════
\section{Costo Variable Unitario ($cv$)}
% ═══════════════════════════════════════════════════════════════════════════════

\begin{enumerate}[label=\textbf{\alph*.}, leftmargin=*, itemsep=0pt]

\item Una empresa de calzado tiene costos variables totales de \$500.000 para una producción de 2.000 pares de zapatos. ¿Cuál es el costo variable unitario?
\answerspace

\item Para fabricar un escritorio, una carpintería gasta \$15.000 en madera y \$3.000 en barniz y clavos. Si estos son sus únicos costos variables, calcule el $cv$ por escritorio.
\answerspace

\item Una fábrica de poleras produce actualmente 5.000 unidades con un costo variable total de \$25.000. Si la producción aumenta a 8.000 unidades, ¿cuál será el nuevo costo variable unitario? \textit{(Pista: el $cv$ se mantiene constante ante cambios en el volumen.)}
\answerspace

\item Una empresa de tecnología gasta \$1.200 en componentes para armar un lote de 40 audífonos. ¿A cuánto asciende el costo variable por cada audífono?
\answerspace

\item El costo total de producir 1.000 unidades es de \$15.000. Si los costos fijos de la empresa son \$5.000, determine primero el costo variable total y luego el costo variable unitario.
\answerspace

\item Una imprenta utiliza 50 gramos de papel especial para cada invitación. Si el gramo de papel cuesta \$0,1, calcule el costo variable unitario por invitación.
\answerspace

\item Una empresa tiene costos fijos de \$10.000 y vende su producto a \$4,5. Si su punto de equilibrio se alcanza al vender 5.000 unidades, despeje el costo variable unitario usando la fórmula $Q = \dfrac{F}{p-c}$.
\answerspace

\item La Empresa A produce 2.000 unidades con un costo variable total de \$10.000. La Empresa B produce 2.500 unidades con un costo variable total de \$15.000. ¿Qué empresa tiene el costo variable unitario más bajo?
\answerspace

\item Un restaurante de pizzas gasta \$2.000 en masa y queso por pizza, y paga una comisión de \$500 al repartidor por cada entrega realizada. ¿Cuál es el costo variable unitario total de una pizza entregada?
\answerspace

\item Una fábrica de jeans tiene costos semanales de producción de \$220 para una producción de 40 jeans. Calcule el costo variable unitario (asumiendo que todos los costos de producción mencionados son variables).
\lastanswerspace

\end{enumerate}

\newpage

% ═══════════════════════════════════════════════════════════════════════════════
\section{Costo Fijo Unitario ($cf$)}
% ═══════════════════════════════════════════════════════════════════════════════

\begin{enumerate}[label=\textbf{\arabic*.}, leftmargin=*, itemsep=0pt]

\item Una empresa paga \$450.000 en costos fijos totales y produce 10.000 unidades. ¿Cuál es su costo fijo unitario?
\answerspace

\item La misma empresa del ejercicio anterior aumenta su producción a 15.000 unidades manteniendo sus costos fijos en \$450.000. ¿Cuál es el nuevo costo fijo unitario?
\answerspace

\item Un taller tiene un arriendo de \$125 y paga \$60 en servicios básicos mensuales. Si produce 100 jeans al mes, calcule el costo fijo unitario considerando ambos gastos como fijos.
\answerspace

\item El costo total de producir 500 lámparas es de \$12.000. Si el costo variable total es de \$7.000, determine primero el costo fijo total y luego el costo fijo unitario.
\answerspace

\item Una fábrica tiene capacidad para 16.000 unidades, pero actualmente solo produce 8.000. Si sus costos fijos totales son \$240.000, calcule el $cf$ actual.
\answerspace

\item Usando los datos del ejercicio anterior, ¿cuál sería el $cf$ si la fábrica decidiera operar a su máxima capacidad (16.000 unidades)?
\answerspace

\item Si el costo fijo unitario de un producto es de \$12 cuando se producen 2.000 unidades, ¿a cuánto ascienden los costos fijos totales de la empresa?
\answerspace

\item Una empresa tiene costos fijos de \$10.000 y alcanza su punto de equilibrio al vender 5.000 unidades. ¿Cuál es el costo fijo unitario en ese punto exacto?
\answerspace

\item La Empresa A tiene costos fijos de \$5.000 y produce 1.000 unidades. La Empresa B tiene costos fijos de \$8.000 y produce 2.000 unidades. ¿Cuál tiene un costo fijo unitario más bajo?
\answerspace

\item Una empresa paga \$3.000 en sueldos administrativos y \$2.000 en seguros de equipos. Si en un mes se fabrican 2.500 unidades, ¿cuál es el costo fijo unitario total?
\lastanswerspace

\end{enumerate}

\newpage

% ═══════════════════════════════════════════════════════════════════════════════
\section{Costo Total ($CT$)}
% ═══════════════════════════════════════════════════════════════════════════════

\begin{enumerate}[label=\textbf{\arabic*.}, leftmargin=*, itemsep=0pt]

\item Una empresa tiene costos fijos mensuales de \$5.000 y un costo variable unitario de \$2. Si produce 1.500 unidades en el mes, ¿cuál es su costo total?
\answerspace

\item Una empresa tiene costos variables de \$200.000 y costos fijos de \$450.000 para una producción de 10.000 unidades. Calcule el costo total de operación.
\answerspace

\item Usando los datos del ejercicio anterior, si la producción aumenta a 15.000 unidades y el costo variable unitario se mantiene en \$20, calcule el nuevo costo total. \textit{(Recuerda: el costo fijo total no cambia.)}
\answerspace

\item Un taller tiene un arriendo de \$125, paga \$60 en servicios y gasta \$237 en materiales para la producción del mes. ¿A cuánto asciende el costo total mensual?
\answerspace

\item Se compra una máquina cuya cuota mensual es de \$308 y el arriendo es de \$80. Para producir 100 unidades, se gastan 20 gramos de tinta por unidad a un precio de \$0,5 el gramo. Calcule el costo total de producir esas 100 unidades.
\answerspace

\item El costo total de producir 2.000 lámparas es de \$18.000 y se sabe que el costo variable unitario es de \$5. Determine a cuánto ascienden los costos fijos de la fábrica.
\answerspace

\item Una empresa alcanza su equilibrio vendiendo 5.000 unidades a un precio de \$4,5 cada una. En el punto de equilibrio, los Ingresos Totales son iguales a los Costos Totales. ¿Cuál es el costo total en ese punto?
\answerspace

\item Un emprendedor paga un seguro fijo de \$1.000 al mes. Cada unidad que fabrica le cuesta \$15 en materias primas y \$5 en mano de obra directa. Calcule el costo total si fabrica 300 unidades.
\answerspace

\item Determine cuál alternativa tiene un Costo Total menor si se planean producir 20 unidades:
  \begin{itemize}
    \item \textbf{Alternativa A:} Costos fijos \$80 y costo variable unitario \$10.
    \item \textbf{Alternativa B:} Costos fijos \$50 y costo variable unitario \$12.
  \end{itemize}
\answerspace

\item Una empresa fabrica 1.000 unidades. Su costo fijo unitario es de \$45 y su costo variable unitario es de \$20. ¿Cuál es el costo total de la producción?
\lastanswerspace

\end{enumerate}

\newpage

% ═══════════════════════════════════════════════════════════════════════════════
\section{Cantidad de Equilibrio ($Q_{eq}$)}
% ═══════════════════════════════════════════════════════════════════════════════

\begin{enumerate}[label=\textbf{\arabic*.}, leftmargin=*, itemsep=0pt]

\item Una empresa tiene costos fijos de \$10.000, costos variables de \$2,5 por unidad y el precio de venta es de \$4,5. Calcule las unidades de equilibrio.
\answerspace

\item Los costos fijos de una empresa son \$5.000 y el precio de venta unitario es \$4,5. Si producir 100 unidades tiene un costo variable de \$250, determine la cantidad de equilibrio. \textit{(Pista: calcula primero el $cv$.)}
\answerspace

\item Usando los datos del ejercicio anterior, calcula el punto de equilibrio expresado en dinero (ventas totales).
\answerspace

\item Una empresa tiene costos fijos de \$17.500. El precio de venta es \$14,5 y el costo variable es de \$180 por cada 175 unidades producidas. Calcule las unidades de equilibrio redondeando al entero más cercano.
\answerspace

\item Una empresa tiene costos fijos de \$8.000, un costo variable unitario de \$10 y vende su producto a \$20. Calcula su $Q_{eq}$. ¿Qué ocurre si el precio de venta sube a \$25?
\answerspace

\item Una fábrica produce a un costo variable de \$5 y vende a \$15. Sus costos fijos suben de \$1.000 a \$1.500 por aumento del arriendo. ¿Cuántas unidades adicionales debe vender para mantenerse en equilibrio?
\answerspace

\item Un emprendedor sabe que su punto de equilibrio es de 500 unidades. Si vende cada unidad a \$100 y su costo variable unitario es de \$60, determina a cuánto ascienden sus costos fijos.
\answerspace

\item Una empresa tiene costos fijos de \$12.000 y un costo variable unitario de \$40. Si su capacidad máxima de producción es de 1.000 unidades y desea que ese sea su punto de equilibrio, ¿a qué precio mínimo debe vender cada producto?
\answerspace

\item Una empresa de jeans tiene costos fijos de \$185 y costos variables de \$15 por unidad. Si el precio de venta por jean es de \$30, calcula el $Q_{eq}$.
\answerspace

\item Determina cuál proyecto requiere vender menos unidades para alcanzar el punto de equilibrio:
  \begin{itemize}
    \item \textbf{Proyecto A:} Fijos \$2.000, Precio \$50, Variable \$30.
    \item \textbf{Proyecto B:} Fijos \$3.000, Precio \$60, Variable \$30.
  \end{itemize}
\lastanswerspace

\end{enumerate}

\newpage

% ═══════════════════════════════════════════════════════════════════════════════
\section{Ingresos Totales ($IT$)}
% ═══════════════════════════════════════════════════════════════════════════════

\begin{enumerate}[label=\textbf{\arabic*.}, leftmargin=*, itemsep=0pt]

\item Una empresa vende su producto a \$12,5. Si logra vender 100 unidades en un mes, ¿cuál es su ingreso total mensual?
\answerspace

\item Una empresa alcanza su punto de equilibrio al vender 5.000 unidades a un precio de \$4,5. ¿A cuánto ascienden sus ingresos totales en ese punto?
\answerspace

\item Se estima una venta mensual constante de \$640. ¿Cuál será el ingreso total acumulado al finalizar los 12 meses?
\answerspace

\item Una empresa produce 10.000 unidades y fija un precio de venta de \$75 por unidad. ¿Cuál sería el ingreso total esperado?
\answerspace

\item Una empresa reporta ingresos totales por \$18.836 y el precio de venta unitario es de \$14,5. ¿Cuántas unidades vendió aproximadamente?
\answerspace

\item Una empresa vende sus productos por \$640 y además vende una máquina usada en \$700. ¿Cuál es el ingreso total percibido en ese periodo?
\answerspace

\item Una empresa vende 2.500 unidades a \$4,5 cada una. Si decide subir el precio a \$6 manteniendo el mismo volumen de ventas, ¿en cuánto aumentan sus ingresos totales?
\answerspace

\item Una empresa tiene capacidad de producción de 16.000 unidades. Si venden toda su capacidad a un precio de \$50, ¿cuál es el ingreso total máximo que pueden generar?
\answerspace

\item Determina cuál proyecto genera mayores ingresos totales:
  \begin{itemize}
    \item \textbf{Proyecto A:} Vende 100 unidades a \$12,5 cada una.
    \item \textbf{Proyecto B:} Vende 80 unidades a \$15 cada una.
  \end{itemize}
\answerspace

\item Una empresa tiene costos fijos de \$17.500 y costos variables de \$1,03 por unidad. Su punto de equilibrio es de 1.299 unidades a un precio de \$14,5. Demuestra que el ingreso total iguala a los costos totales en ese punto.
\lastanswerspace

\end{enumerate}

\newpage

% ═══════════════════════════════════════════════════════════════════════════════
\section{Punto de Equilibrio}
% ═══════════════════════════════════════════════════════════════════════════════

\begin{enumerate}[label=\textbf{\arabic*.}, leftmargin=*, itemsep=0pt]

\item Una empresa tiene costos fijos de \$10.000, un precio de venta de \$4,5 y costos variables de \$2,5 por unidad. ¿Cuántas unidades debe vender para alcanzar el punto de equilibrio?
\answerspace

\item Los costos fijos son de \$5.000 y el precio de venta unitario es de \$4,5. Si producir 100 unidades tiene un costo variable total de \$250, determina la cantidad de equilibrio.
\answerspace

\item Con los datos del ejercicio anterior ($Q_{eq} = 2.500$ unidades y precio \$4,5), calcula el punto de equilibrio expresado en unidades monetarias.
\answerspace

\item Una empresa tiene costos fijos de \$17.500. El precio de venta es \$14,5 y el costo variable es de \$180 por cada 175 unidades. Calcula las unidades de equilibrio redondeando al entero más cercano.
\answerspace

\item Con los datos del ejercicio anterior ($Q_{eq} = 1.299$ unidades y precio \$14,5), determina el ingreso total necesario para no tener pérdidas.
\answerspace

\item Una empresa tiene costos fijos de \$450.000 y un costo variable unitario de \$20. Si fijan un precio de venta de \$70, ¿alcanza su punto de equilibrio dentro de su capacidad máxima de 16.000 unidades?
\answerspace

\item Una empresa de jeans tiene costos fijos de \$185, precio de venta \$30 y costo variable \$15. Calcula el $Q_{eq}$. ¿Qué ocurre si los costos fijos suben a \$300?
\answerspace

\item Un emprendedor tiene costos fijos de \$10.000 y costos variables de \$2,5 por unidad. Si solo puede producir y vender 2.000 unidades, ¿a qué precio debe vender cada una para que ese sea su punto de equilibrio?
\answerspace

\item Si una empresa logra bajar su costo variable unitario de \$2,5 a \$1,5, manteniendo costos fijos de \$5.000 y un precio de \$4,5, ¿en cuántas unidades se reduce su punto de equilibrio?
\answerspace

\item Demuestra que en el punto de equilibrio ($Q = 5.000$, $p = 4{,}5$, $c = 2{,}5$, $F = 10.000$) los Ingresos Totales son exactamente iguales a los Costos Totales, es decir que $p \cdot Q = F + c \cdot Q$.
\lastanswerspace

\end{enumerate}

\newpage

% ═══════════════════════════════════════════════════════════════════════════════
\section{Valor Presente ($VP$)}
% ═══════════════════════════════════════════════════════════════════════════════

\begin{enumerate}[label=\textbf{\arabic*.}, leftmargin=*, itemsep=0pt]

\item Si una inversión promete entregarte \$1.200 en 5 años, ¿cuánto vale ese dinero hoy si la tasa de interés es del 3\%?
\answerspace

\item Calcula el valor hoy de esos mismos \$1.200, pero considerando que se recibirán en 3 años en lugar de 5, manteniendo la tasa del 3\%.
\answerspace

\item Patricio Lucas tiene un contrato por 7 temporadas donde recibirá \$400.000 por cada una. Si su costo alternativo es del 10\%, ¿cuál es el valor presente total de ese contrato?
\answerspace

\item Siguiendo el caso anterior, si le ofrecen \$1.900.000 pagados hoy mismo en lugar de las cuotas anuales, ¿cuál opción tiene un valor presente mayor?
\answerspace

\item Si tienes \$100 hoy y la tasa de interés es del 10\%, ¿cuánto valdrían en un año si los depositas en el banco? ¿Cuál es el valor presente de esos \$110 futuros?
\answerspace

\item Una máquina se vende en \$700 en el mes 12 de un proyecto. Calcula el valor presente de esa venta con una tasa mensual del 1\%.
\answerspace

\item Una empresa necesita tener \$10.000 en 6 años para renovar equipos. Con una tasa del 6,5\%, ¿cuánto debe invertir hoy para alcanzar esa meta?
\answerspace

\item Determina el VP de un pago de \$500 que se recibirá en 2 años, utilizando una tasa de rentabilidad del 4\%.
\answerspace

\item Un proyecto espera vender sus máquinas por \$400 al finalizar el segundo año. Calcula el VP de este ingreso si la tasa exigida es del 10\%.
\answerspace

\item Calcula el VP de \$1.000 a recibir en 1 año con un 10\% de tasa. ¿Por qué ese resultado confirma que un peso futuro vale menos que un peso hoy?
\lastanswerspace

\end{enumerate}

\newpage

% ═══════════════════════════════════════════════════════════════════════════════
\section{Valor Presente Neto ($VAN$)}
% ═══════════════════════════════════════════════════════════════════════════════

\begin{enumerate}[label=\textbf{\arabic*.}, leftmargin=*, itemsep=0pt]

\item Para una inversión inicial de \$1.000 se proyectan 4 cuotas de retorno de \$350 cada una. Si el inversionista exige un 10\% de rentabilidad, ¿cuál es el VAN?
\answerspace

\item Basado en el resultado anterior ($VAN = \$109{,}45$), ¿le aconsejarías al inversionista realizar el proyecto? Justifica tu respuesta.
\answerspace

\item Calcula el VAN de un proyecto con inversión inicial de \$25.000 y los siguientes flujos: Año 1: \$8.000, Año 2: \$6.000, Años 3 y 4: \$0, Año 5: \$10.000, Año 6: \$10.000. Usa una tasa del 6,5\%.
\answerspace

\item Un emprendedor con costo de oportunidad del 8\% debe elegir entre:
  \begin{itemize}
    \item \textbf{Alt.\ 1:} Inversión de \$1.200 con retornos de \$354, \$123 y \$1.876 en los meses 1, 2 y 3.
    \item \textbf{Alt.\ 2:} Arrendar una máquina con flujos de \$845, \$656 y \$754.
  \end{itemize}
  ¿Qué alternativa es más conveniente según el VAN?
\answerspace

\item Se compra una máquina de \$1.500 que produce 100 unidades mensuales vendidas a \$12,5 cada una. Si el interés es del 0,4\% mensual, calcula el VAN al finalizar el periodo.
\answerspace

\item Calcula el VAN de una inversión de \$2.000 que solo devuelve \$500 anuales por 3 años con una tasa del 15\%. ¿Es rentable?
\answerspace

\item Si un proyecto tiene un VAN de 0, ¿qué significa esto para el inversionista respecto a su tasa mínima exigida?
\answerspace

\item Un proyecto de diseño requiere \$5.000 de inversión y genera \$2.000 anuales durante 3 años. Si la tasa es del 7\%, calcula el VAN.
\answerspace

\item Calcula el VAN de un proyecto que requiere \$1.500, genera \$600 anuales por 3 años y permite vender el equipo al final en \$400, con una tasa del 10\%.
\answerspace

\item Un inversionista evalúa dos proyectos con tasa del 10\%:
  \begin{itemize}
    \item \textbf{Proyecto X:} Inversión \$3.000; flujos anuales de \$1.200 por 3 años.
    \item \textbf{Proyecto Y:} Inversión \$4.000; flujos anuales de \$1.500 por 3 años.
  \end{itemize}
  Calcula el VAN de ambos y determina cuál es más conveniente.
\lastanswerspace

\end{enumerate}

\vspace{1em}
\noindent\rule{\textwidth}{0.8pt}
\begin{center}
  \small\textit{Recuerda: VAN $> 0$ → proyecto conveniente \quad · \quad VAN $= 0$ → indiferente \quad · \quad VAN $< 0$ → no conveniente}
\end{center}

\newpage

% ═══════════════════════════════════════════════════════════════════════════════
\section{ITEM-01: Diferenciación de Costos Fijos y Variables}
% ═══════════════════════════════════════════════════════════════════════════════

\noindent\textit{Contexto: Un nuevo emprendimiento de fabricación de jeans está organizando sus cuentas mensuales.}

\vspace{0.5em}
\noindent Clasifica los siguientes desembolsos como \textbf{Costo Fijo (CF)} o \textbf{Costo Variable (CV)} y justifica tu respuesta:

\begin{enumerate}[label=\textbf{\arabic*.}, leftmargin=*, itemsep=0pt]

\item El pago del alquiler del taller de costura.
\answerspace

\item La compra de tela de mezclilla (materia prima) para los pantalones.
\answerspace

\item El seguro contra incendios de la maquinaria.
\answerspace

\item Las comisiones pagadas a los vendedores por cada jean vendido.
\answerspace

\item El servicio de internet de banda ancha de la oficina.
\lastanswerspace

\end{enumerate}

\newpage

% ═══════════════════════════════════════════════════════════════════════════════
\section{ITEM-02: El Efecto del Volumen en los Costos Unitarios}
% ═══════════════════════════════════════════════════════════════════════════════

\noindent\textit{Contexto: La empresa ``Los Pitufos'' produce actualmente 10.000 unidades. Sus costos totales son de \$450.000 en costos fijos y \$200.000 en costos variables.}

\vspace{0.5em}

\begin{enumerate}[label=\textbf{\arabic*.}, leftmargin=*, itemsep=0pt]

\item Calcula el \textbf{costo fijo unitario} y el \textbf{costo variable unitario} para el nivel de producción actual (10.000 unidades).
\answerspace

\item \textbf{Escenario de crecimiento:} Si la producción aumenta a 15.000 unidades, ¿cuál será el nuevo costo fijo unitario y el nuevo costo variable total?
\lastanswerspace

\end{enumerate}

\newpage

% ═══════════════════════════════════════════════════════════════════════════════
\section{ITEM-03: Cálculo del Punto de Equilibrio}
% ═══════════════════════════════════════════════════════════════════════════════

\noindent\textit{Contexto: Un restaurante analiza la rentabilidad de su plato estrella, ``Carne con arroz'', con los siguientes datos:}

\vspace{0.4em}
\begin{itemize}[leftmargin=2cm, itemsep=2pt]
  \item Precio de venta por plato: \$14,5
  \item Costos variables (ingredientes y empaques): \$2,5 por plato
  \item Costos fijos mensuales (alquiler, sueldos, servicios): \$10.000
\end{itemize}
\vspace{0.4em}

\begin{enumerate}[label=\textbf{\arabic*.}, leftmargin=*, itemsep=0pt]

\item Calcula cuántos platos debe vender el restaurante al mes para alcanzar su punto de equilibrio (beneficio = 0).
\answerspace

\item Si el restaurante logra vender 1.200 platos, ¿está generando ganancias o pérdidas? Justifica usando la fórmula de Ingresos Totales vs.\ Costos Totales.
\lastanswerspace

\end{enumerate}

\newpage

% ═══════════════════════════════════════════════════════════════════════════════
\section{ITEM-04: Flujo de Caja}
% ═══════════════════════════════════════════════════════════════════════════════

\noindent\textit{Contexto: Tienes un emprendimiento que produce y vende jeans. Construye un flujo de caja mensual para el primer año con los siguientes datos:}

\vspace{0.4em}
\begin{enumerate}[label=\textbf{\arabic*.}, leftmargin=2cm, itemsep=3pt]
  \item \textbf{Ventas:} \$640 cada mes.
  \item \textbf{Inversión en Maquinaria:}
    \begin{itemize}
      \item Mes 1: se compra una máquina por \$1.850, pagada en 6 cuotas iguales de \$308 (enero a junio).
      \item Mes 12: se vende esa máquina por \$700 y se compra una nueva por \$2.000.
    \end{itemize}
  \item \textbf{Costos Variables:} \$55 a la semana $\times$ 4,3 semanas = \$237 al mes.
  \item \textbf{Costos Fijos mensuales:} Alquiler \$125 · Servicios \$60.
\end{enumerate}
\vspace{0.6em}

\noindent Completa la tabla de flujo de caja a continuación:

\vspace{0.8em}
\noindent\resizebox{\textwidth}{!}{%
\renewcommand{\arraystretch}{2}
\begin{tabular}{|l|*{12}{c|}}
\hline
\textbf{Concepto} & \textbf{M1} & \textbf{M2} & \textbf{M3} & \textbf{M4} & \textbf{M5} & \textbf{M6} & \textbf{M7} & \textbf{M8} & \textbf{M9} & \textbf{M10} & \textbf{M11} & \textbf{M12} \\
\hline
Ingresos por ventas & & & & & & & & & & & & \\
\hline
Venta de máquina & & & & & & & & & & & & \\
\hline
Cuota máquina & & & & & & & & & & & & \\
\hline
Compra nueva máquina & & & & & & & & & & & & \\
\hline
Costos variables & & & & & & & & & & & & \\
\hline
Alquiler & & & & & & & & & & & & \\
\hline
Servicios & & & & & & & & & & & & \\
\hline
\textbf{Flujo Neto} & & & & & & & & & & & & \\
\hline
\end{tabular}%
}

\vspace{1.5cm}

\newpage

% ═══════════════════════════════════════════════════════════════════════════════
\section{ITEM-05: Valor del Dinero en el Tiempo}
% ═══════════════════════════════════════════════════════════════════════════════

\begin{enumerate}[label=\textbf{\arabic*.}, leftmargin=*, itemsep=0pt]

\item \textbf{Valor Presente (VP):} Una persona recibirá \$1.000.000 dentro de 3 años. La tasa de interés anual es del 10\%. ¿Cuál es el valor presente de ese dinero hoy?

\vspace{0.3em}
\noindent\textit{Datos: $VF = 1.000.000$ \quad $i = 0{,}10$ \quad $n = 3$}
\answerspace

\item \textbf{Valor Futuro (VF):} Inviertes hoy \$500.000 a una tasa de interés anual del 8\% durante 4 años. ¿Cuánto dinero tendrás al final?

\vspace{0.3em}
\noindent\textit{Datos: $VP = 500.000$ \quad $i = 0{,}08$ \quad $n = 4$}
\answerspace

\item \textbf{Comparación de contratos (VP):} A un cantante le ofrecen dos tipos de contrato. Analiza cuál le conviene más:
  \begin{itemize}
    \item \textbf{Opción A:} US\$~300.000 al final de cada año durante 6 años.
    \item \textbf{Opción B:} Un pago único de US\$~1.300.000 hoy.
  \end{itemize}
  La tasa de descuento (costo de oportunidad) es del 9\% anual. Calcula el VP de la Opción A y compáralo con la Opción B.
\answerspace

\item \textbf{Valor Actual Neto (VAN):} Calcula el VAN del siguiente proyecto si la tasa de interés anual es del 25\%:

\vspace{0.5em}
\begin{center}
\renewcommand{\arraystretch}{1.4}
\begin{tabular}{|c|r|}
\hline
\textbf{Año} & \textbf{Flujo} \\
\hline
0 & $-$300.000 \\
1 & 145.000 \\
2 & 95.000 \\
3 & 37.000 \\
4 & 256.000 \\
\hline
\end{tabular}
\end{center}
\vspace{0.4em}
¿Es conveniente realizar este proyecto? Justifica.
\lastanswerspace

\end{enumerate}

\vspace{1em}
\noindent\rule{\textwidth}{0.8pt}
\begin{center}
  \small\textit{Recuerda: VAN $> 0$ → proyecto conveniente \quad · \quad VAN $= 0$ → indiferente \quad · \quad VAN $< 0$ → no conveniente}
\end{center}

\end{document}